\scanpage{5}
Noter que $\mathrm{cd}_{p}^{\text{ét}}(k)$ est égal à $0$ ou $1$, suivant que \deleted{toute} extension \inserted{finie} étale de $k$ est d'ordre premier à $p$, ou non.

Si $G$ est de t. m., alors $G$ est dual d'un groupe fini étale. Ce dernier est quotient d'un module galoisien, \unsure{libre} de rang fini sur $\mathbb{Z}$, donc $G$ est contenu dans un tore, \unsure{donc on a} une suite exacte
\[
0\longrightarrow G\longrightarrow T\longrightarrow T'\longrightarrow 0
\]
$T$ et $T'$ des tores. D'où le

\uline{Lemme} 2 \quad Pour que $\mathrm{cd}_{p}^{\text{t.m.}}(k)$ soit $\leq n$, il f. et s. que les conditions suivantes soient vérifiées : pour tt tore $T$ sur $k$, $H^{i}(k,T)(p)$ est nul pour $i>n$, et $p$-divisible si $i=n$. [Ou encore : pour tte extension finie \unsure{séparable} $k'$ de $k$ [il suffit d'ordre premier à $p$] \deleted{$\mathrm{Br}(k)(p)$} $H^{i}(k',G_{m})=0$
% margin note written sideways; paired loop carries on the square bracket.
\margin{si $i>n$, et $H^{n}(k',G_{m})$ est $p$-divisible]}

Comme on a $0\longmapsto {}_{p}T\longrightarrow T\xrightarrow{p}T\longrightarrow 0$

d'où
\[
\deleted{\left\{ {}_{p}H^{\unsure{i|\prime}}(T)=\operatorname{Im} H^{i}({}_{p}T)\right.}
\]

% archive sheet number: 4

